% 原书第 17 页（PDF 第 30 页）
\BookPage{17}{30}
\BookHeading{清代中期扬州市镇经济水平和民众生活初探—以刑科题本档案资料为基础}{清代中期扬州市镇经济水平和民众生活初探\\[3pt]—以刑科题本档案资料为基础}{王跃生*}
清代中期，扬州是东南地区的一大都会，其繁盛情景当时多种文献中有所记载。但在我看来，官私文献关于普通民众生存方式的记述有线条较粗之感（有关扬州盐商、文人士大夫生活的史料甚多，而对一般民众则涉及较少），可以说，我们对当时扬州城镇经济水平的认识仍是比较模糊的。最近在查阅档案资料时，我发现，档案中有一些个案对当事者个人生存方式有所反映。而多件个案汇集起来，一定程度上将可揭示出特定区域经济发展水平和民众生活状况。个案讯息尽管不能为认识一个地区经济发展水平和民众生活提供全景式视野，却也不失为一个窗口。本文将以乾隆朝、嘉庆朝和道光朝初期的刑科题本（土地债务类、婚姻家庭类）档案为基础，以扬州所属江都县和甘泉县（这两地应为当时扬州的核心区域）个案为观察对象，从局部和某一侧面认识清代中期扬州的城镇经济发展水平，进而对当时城镇居民的生存方式有所了解。

\BookSubheading{一、扬州的商业服务水平}

清朝中期，扬州是一个以消费为主的城市，商业服务相对发达。这一点档案中有所表现。

\BookSubheading{（一）城镇居民日常生活与市场联系密切}

对传统时代民众日常生活市场化程度的认识，吃、住、住、行、用所需物品来源是重要的考察视角。

\vfill
\begin{BookNotes}
\BookNote{*}{作者任职于中国社会科学院人口与劳动经济研究所。}
\end{BookNotes}

% 原书第 18 页（PDF 第 31 页）
\BookPage{18}{31}
我们从档案中看到想像不到的扬州人生活场景：居民在家洗澡时所用热水从市场上购买。

\begin{BookQuote}
蔡伍供：江都县人，年三十三岁。小的与李宏素识无嫌。嘉庆二十三年三月十二日，李宏租住小的房屋两间，每月租钱三百文，当收押租钱八千文。七月十五日，李宏要另租住房，叫小的退还押租。小的原应许他八月十一日交还，至期未经措交。那夜一更多天，小的因洗澡，提了木桶到茶炉上买得热水回来。李宏坐在门外小板凳上吃酒，看见小的走至，向小的催逼押租，小的恳缓。那时李宏已醉，开口就骂。小的生气，提起水桶吓说泼水烫他。不料李宏陡然站起，头顶撞翻水桶，以致碰伤额门，并热水烫伤他周身，…。1
\end{BookQuote}

蔡五在所欠押租未还之际仍从茶房买热水洗澡，表明洗热水澡已成为他的生活方式，是生活比较讲究的体现。更重要的是，茶炉在供应茶水的同时，还兼卖洗澡用热水。可谓有需求就有供给。当然，我们相信，热水的价格也不会很高，否则居民宁愿自己烧水。一般来说，茶水店烧开水具有规模优势，其成本将低于家庭自己烧水费用。这件小事反映出清代中期扬州城市商业服务的细致。

而扬州更有专门的澡堂一时称「混堂」为居民提供服务。

\begin{BookQuote}
柳在友供：江都县人，柳兆昌是小的哥子，开混堂生理。2
\end{BookQuote}

值得注意的是，这些澡堂烧水所用柴草也有专门的草行提供。

此外，扬州还有专门的洗衣店。

\begin{BookQuote}
杨四供：江都人，二十九岁，父母已故，并无兄弟妻子。向撑木排度日。乾隆五十四年四月九日，小的有一衣交陈俞氏擦洗，许钱七文。3
\end{BookQuote}

此种服务业的存在为未婚、没有家室的男性提供了生活便利。

\vfill
\begin{BookNotes}
\BookNote{1.}{杜家骥主编，《清嘉庆朝刑科题本社会史料辑刊》（以下简称《史料辑刊》）第2册，天津古籍出版社，2008年，第749页。}
\BookNote{2.}{中国第一历史档案馆，刑科题本，土地债务类，乾隆五十六年九月二十三日，江宁巡抚长麟题。}
\BookNote{3.}{中国第一历史档案馆，刑科题本，土地债务类，乾隆五十年十月十九日，江宁巡抚闵鹗元题。}
\end{BookNotes}

% 原书第 19 页（PDF 第 32 页）
\BookPage{19}{32}
\BookSubheading{（二）物流环节比较健全}

商业物流体现为同一领域商业经营者之间的分工和合作。只有健全的物流服务才能促使商业活动提升水平。个案显示，清代中期扬州的物流环节是比较健全的。

开杂货店者为坐商，有行商将货物送来。不少个案显示，坐商实际通过赊账来收购货物，等货物买出后再与供货商人结算。

这种经营形式表明，扬州地区的商业网络已经形成，商业信用具有基础。即使同一行业的商业经营，也有多个环节相互配套。它是商业分工程度较高的表现。

酒行与酒供应商之间的关系：

\begin{BookQuote}
张能供：江都县人，四十三岁，小的开酒行生理。何云存代槽坊挑酒贩卖，小的常与他交易。小的经手各酒店处，共欠何云存三千零四十文钱，除他陆续收过外，还欠他一百文。4
\end{BookQuote}

开酒店者所用酒不必自己派人去探购，而由专门推销者供应。

酒店与肉店之间的关系：

\begin{BookQuote}
袁文通供：甘泉县人，二十六岁，小的开酒馆营生，尹成贤是开肉店的，住居邻近，一向和好。小的向与他买肉，陆续尾欠，积有三十两银子没有还他。5
\end{BookQuote}

在我看来，物流服务健全的直接作用是减少了开店者的用工成本，少雇或不雇人手即可运转，开办商店所需资金门槛因此而降低，它对市镇商业的发展具有推动作用。

当然，经营过程中的赊账行为有利也有弊。对酒店经营者来说，赊账可以解决资金短缺问题。但赊欠到较高的数额、较长的时间仍不能结算，则会失去信用，使供应链或合作中断。同时对被赊者来说，资金短缺状况则会出现，经营的困难程度将提高，严重者则可能关门歇业。正因为有这种忧虑，经营者往往不得不找欠钱之人反复索欠。

\vfill
\begin{BookNotes}
\BookNote{4.}{中国第一历史档案馆，刑科题本，土地债务类，乾隆二十七年三月初一日，江苏巡抚陈宏谋题。}
\BookNote{5.}{中国第一历史档案馆，刑科题本，土地债务类，乾隆四年九月十五日，江宁巡抚张渠题。}
\end{BookNotes}

% 原书第 20 页（PDF 第 33 页）
\BookPage{20}{33}
\BookSubheading{（三）商业活动形成规模，管理比较规范}

在扬州，民众开店，特别是开办规模较大的经营行业，需向当地政府申请。它表明官方把基本商业活动纳入管理体系中，这有助于建立和保持地方商业秩序。

\begin{BookQuote}
王德供：甘泉县人，领帖开张鲜果行生理。6
\end{BookQuote}

这里的「领帖」实际是从官方机构获得开业执照，是正规的「坐商」。它意味着当地商业店面较多，形成了商业规模。而同案之中，另一涉案人的供述说明了这一点。

\begin{BookQuote}
曹大供：山东人，嘉庆十八年三月初三日，小的从原籍贩梨七十七包来到扬州（江都县一笔者注）鲜果行内售卖。同乡霍泳贵、霍泳泰走来劝小的运二、三十包赴别行分售，以免货多阻滞。7
\end{BookQuote}

可见，扬州当地鲜果行较多。

该案例表明，外地客商将水果运来，可以利用鲜果行所提供的场地，自行销售。果行根据销售数量收取管理费用。它与现代城市中所建立的蔬菜、水果交易市场很相像。

下面的案例则进一步告诉我们，领帖经营并非对个别商行的要求。

\begin{BookQuote}
许应陇供：甘泉县人，五十四岁。小的与杨德淮都领帖开菜行生理。嘉庆二十一年上，小的没钱使用，把菜行地段立契绝卖与杨德淮，得价十五千文。8菜行地段可能是大的菜市场中的一个永久摊位，归个人所有。
\end{BookQuote}

在我看来，「领帖」从事经营活动是城市商业管理规范化的一种体现。

当然，也有地点不固定的零售商贩，他们很可能不必申领执照。

\begin{BookQuote}
陈三供：甘泉县人，六十二岁。乾隆五十七年六月二十三日，小的拿了一篮西
\end{BookQuote}

\vfill
\begin{BookNotes}
\BookNote{6.}{中国第一历史档案馆，刑科题本，土地债务类，嘉庆十八年十一月初一日，江宁巡抚朱理题。}
\BookNote{7.}{中国第一历史档案馆，刑科题本，土地债务类，嘉庆十八年十一月初一日，江宁巡抚朱理题。}
\BookNote{8.}{中国第一历史档案馆，刑科题本，土地债务类，道光五年十月二十九日，江苏巡抚陶澍题。}
\end{BookNotes}

% 原书第 21 页（PDF 第 34 页）
\BookPage{21}{34}
\noindent 瓜歇在陈有饭店外棚下零卖。9

孔文供：甘泉县人，五十四岁，小的与李士一都住在陈国泰酱坊相近。李士一向在陈国泰酱坊门首摆摊子卖炒熟蚕豆等杂货，小的是挑担卖荸荠等水果的。10

我们认为，扬州的商业具有规范管理和灵活经营相结合的特征。

\BookSubheading{（四）商业经营的专业化程度高}

\noindent 扬州的商店和专项经营者多以「行」相称。不仅有粮行、鲜果行，而且还有船行、草行（提供燃料）。「行」体现了经营的专业性和专门性。

王世太供：甘泉县人，小的平日开船行生理，与船户何起受相好。乾隆四十一年二月十五日，何起受向小的借去七百钱，约俟揽载清还。11

张玉供：江都县人，挑脚营生。乾隆五十六年二月二十五日，草行奚恒雇小的同刘老汉挑草两担，送到柳兆昌店内（开澡堂店）。12

开澡堂者烧水所需燃料有专营柴草者负责供应。

\BookSubheading{(五)商业竞争有一定表现}

\noindent 同行之间竞争客户是自由市场经济下常见的现象。这在扬州的多种行业表现出来。

豆腐店之间经营之争。徐士友供：小的是江都县人，同哥子徐士英开张腐店，与虞臣松邻居无嫌，只因虞臣松近日也卖豆腐，他把小的家主顾都揽去了。乾隆三十八年八月十七日早上，小的气他不过，走到他家把他磨子拿了一片回来，原说大家卖不成豆

\vfill
\begin{BookNotes}
\BookNote{9.}{中国第一历史档案馆，刑科题本，土地债务类，乾隆五十八年二月二十二日，大学士管理刑部事务阿桂题。}
\BookNote{10.}{中国第一历史档案馆，刑科题本，土地债务类，乾隆十二年三月二十三日，刑部尚书阿克敦题。}
\BookNote{11.}{中国第一历史档案馆，刑科题本，土地债务类，乾隆四十二年三月初六日，大学士管理刑部事务舒赫德题。}
\BookNote{12.}{中国第一历史档案馆，刑科题本，土地债务类，乾隆五十六年九月二十三日，江宁巡抚长麟题。}
\end{BookNotes}

% 原书第 22 页（PDF 第 35 页）
\BookPage{22}{35}
\begin{BookQuote}
腐。虞臣松回来向小的索讨，…。13
\end{BookQuote}

这种竞争对消费者是有利的，由此他们能更自由地选择购买对象。

船运者之间竞争。通过压价招揽客户。

\begin{BookQuote}
徐立成供：小的向有划船一只，雇陈五在外揽载，与刘八认识。嘉庆十七年十月初八日晚陈五告诉小的说，船在镇江地方见有客货三担，刘八的儿子刘大金子先去揽载，因较量船价，未成；客人自来雇他船支，他就减价装送的话。初九日，小的在街扫地，刘八走来说，他儿子刘大金子在镇江揽了三担客货，被陈五减价揽去装载，斥说陈五不是，并向小的理论…。14
\end{BookQuote}

这些案例显示出，清代中期，扬州的城镇市场非常活跃且形成较完善的商业服务体系；扬州民众的经营、就业和消费行为与市场有比较密切的联系。

\BookSubheading{二、扬州人的职业类型}

在农业经营为主的环境中，人们的就业和谋生方式是比较简单的。而作为以非农人口为主的市镇，特别是像扬州这样的商业相对发达地区，民众从事的职业又有什么特征？从档案中我们看到扬州人的职业类型相对集中几个方面。

\BookSubheading{（一）以船运为业者}

1. 绎夫

运河是扬州地区主要水上运输通道，船则为主要交通工具，拉绎成为当地一个重要职业。

陈九供：江都县人，三十九岁，佣趁度日。嘉庆十九年八月十五日二更时小的拉綷回到宝应汜水镇地方，适于德贵帮撑，周围米船停泊在彼，要雇小的拉到清

\vfill
\begin{BookNotes}
\BookNote{13.}{中国第一历史档案馆，刑科题本，土地债务类，乾隆三十九年二月十五日，江宁巡抚萨载题。}
\BookNote{14.}{中国第一历史档案馆，刑科题本，土地债务类，嘉庆十八年六月初七日，江宁巡抚朱理题。}
\end{BookNotes}

% 原书第 23 页（PDF 第 36 页）
\BookPage{23}{36}
\noindent 江，需钱二百文15。2. 水手绎夫是从岸上拉船，船上也需有摇船之人，后者称为水手。水手也是扬州的一个重要职业。

高松山供：甘泉县人，二十四岁。乾隆四十八年十二月，小的央丁有学荐在莫魁英船上做水手，讲定每月工钱二千五百文。16同案中莫魁英父亲莫乔松对儿子的职业如此说明：莫魁英为小的长子，摇船生理。可见水手就是摇船之人。3. 船运主船运主中既有靠雇水手经营者，也有以自己之船直接营运客货为生者。徐立成供：高邮州人，寄居江都。小的向有划船一只，雇陈五在外揽载。17王太仓供：江都人，娶有妻子，驾船营生，父故母存。母夏氏随船度日。18王喜子供：小的江都人，二十八岁，平日撑木排营生。19以船为业者不仅有男性，还有女性。

曹氏供：小妇人甘泉县邵伯镇人，十八岁。父亲已死两年，母亲与小妇人并妹子驾船度日。20

\vfill
\begin{BookNotes}
\BookNote{15.}{中国第一历史档案馆，刑科题本，土地债务类，嘉庆二十年九月二十三日，大学士管理刑部事务董诰题。}
\BookNote{16.}{中国第一历史档案馆，刑科题本，土地债务类，乾隆四十九年（此件档案月份和日期不全一笔者注)，江宁巡抚闵鹗元题。}
\BookNote{17.}{中国第一历史档案馆，刑科题本，土地债务类，嘉庆十八年六月初七日，江宁巡抚朱理题。}
\BookNote{18.}{中国第一历史档案馆，刑科题本，土地债务类，道光二年十一月二十二六日，大学士管理刑部事务戴均元题。}
\BookNote{19.}{中国第一历史档案馆，刑科题本，土地债务类，乾隆五十年十月十九日，江宁巡抚闵鹗元题。}
\BookNote{20.}{中国第一历史档案馆，刑科题本，婚姻家庭类，乾隆三十一年八月初八日，江苏巡抚明德题。}
\end{BookNotes}

% 原书第 24 页（PDF 第 37 页）
\BookPage{24}{37}
\BookSubheading{（二）以开客栈为业者}

\noindent 扬州作为流动人口相对较多的地区，开客栈成为一个重要职业。既然拉成为一种职业，而它又是流动性职业，有很大可能带动住宿行业的发展。我们看到，扬州一带有专门供绎夫居住的客栈。开饭店多为食宿一体经营。

沈车子供：父亲沈朝观开饭店生理，盐城人邰子良无船可撑船时在此吃饭住宿。21有的所开客栈直接叫客寓。

高春生供：小的开张客寓生理。22

\BookSubheading{（三）以售卖各色货物为业}

\noindent 这一行业是扬州人从业的主流，其中又包含许多类型。1. 以卖肉为业者

卜广子供：小的江都人，开肉店生理。嘉庆五年二月小的向吴泰万赊猪二口，讲定价钱四千八百文。23

李廷佐供：小的江都人，李廷位是小的兄弟，开张猪行生理。

同案蒋子珍供：四十三岁，李廷位是开张猪行的，小的常到他行里买猪。24买者主要目的是宰杀生猪，销售猪肉。不少肉店都雇有帮工者：

\vfill
\begin{BookNotes}
\BookNote{21.}{中国第一历史档案馆，刑科题本，土地债务类，乾隆二十年一月二十日，署理两江总督尹继善题。}
\BookNote{22.}{中国第一历史档案馆，刑科题本，土地债务类，道光元年五月十六日，江宁巡抚魏元煜题。}
\BookNote{23.}{中国第一历史档案馆，刑科题本，土地债务类，嘉庆六年四月二十三日，大学士管理刑部事务董诰题。}
\BookNote{24.}{中国第一历史档案馆，刑科题本，土地债务类，乾隆二年九月二十三日，大学士兼管刑部尚书事务徐本题。}
\end{BookNotes}

% 原书第 25 页（PDF 第 38 页）
\BookPage{25}{38}
\begin{BookQuote}
孔登魁供：小的江都人，向在卜广子肉店帮伙。25
\end{BookQuote}

2.以卖杂货为业者

\begin{BookQuote}
胡有兴供：江都县人，年三十六岁，开杂货店生理。26
\end{BookQuote}

3.以开茶炉卖茶水者

开茶炉卖茶水个案前面已有提及。

另一个案中，丈夫和长子出外佣工，主妇与小儿子在家，并开茶水店。

\begin{BookQuote}
万田氏供：四十五岁，江都县人，丈夫万松茂同大儿子万荣增在外佣工。小妇人在门首开茶店生理，常叫吕美挑水送来。27
\end{BookQuote}

这一个案实际又引出另一职业一为茶店送水者。同时它也再次印证了当地物流服务的完备状况：有人专门为茶水店供应凉水。更进一步，开茶水店者将水烧开后，不仅卖茶，而且会卖开水。

4.以开酒店为业者

当时的酒店实际主要是小饭馆或小酒馆，提供酒和基本食物，它与居民生活关系密切。

江都县民严春源欠族人严春荣钱，严春荣将其棉袄拿去，到酒店里押酒吃。严春源不得不到酒店替他付酒钱，然后把棉袄取回28。可见，酒店客人不仅以现钱消费，而且可以通过赊账和押物品来消费。

但赊账会埋下纷争的种子。

\begin{BookQuote}
林春华供：江都县人，开酒店生理，张玉曾在小的店内陆续赊欠酒钱
\end{BookQuote}

\vfill
\begin{BookNotes}
\BookNote{25.}{中国第一历史档案馆，刑科题本，土地债务类，嘉庆六年四月二十三日，大学士管理刑部事务董诰题。}
\BookNote{26.}{中国第一历史档案馆，刑科题本，土地债务类，嘉庆九年九月二十三日，江苏巡抚汪志伊题。}
\BookNote{27.}{中国第一历史档案馆，刑科题本，婚姻家庭类，嘉庆十二年三月初三日，江宁巡抚汪日章题。}
\BookNote{28.}{《史料辑刊》，第1册，第307页。}
\end{BookNotes}

% 原书第 26 页（PDF 第 39 页）
\BookPage{26}{39}
\noindent 二百五十八文，屡次向讨不还。29许二供：江都县人，开酒店生理，熟识的戴启元陆续赊酒欠，该小的五钱八分七厘酒钱，小的时常向他索讨。30也有酒店欠客户的钱。前面个案中酒店欠肉店商人、供酒商人的钱即如此。5. 以开磨坊为业者扬州豆腐的优良品质远近闻名，从业者不在少数。孙应成供：江都县人，与妻子苏氏开豆腐店生理。31而不少豆腐店雇人经营。高秀供：江都县人，小的开张磨坊。田三雇在小的家佣工有十六年了。32可见，雇主和雇工在这一行的就业比较稳定，非短期行为。同案庄三供：小的是丹徒县人，凭谢文陇荐到高秀磨坊里做工。33 6.以经营粮食收购和销售为业者当地有比较正规的「粮行」业务。7.开各种面食店者宋其祥供：六十四岁，向开点心店生理。张大先在小的店内帮伙，后来辞出，他欠少小的钱一千五百文文没还。同案张大供：小的前在宋其祥店内帮伙，欠他钱一千五百文没还，后辞出，又到蔡家沅店内帮伙。34

\vfill
\begin{BookNotes}
\BookNote{29.}{中国第一历史档案馆，刑科题本，土地债务类，乾隆四十三年二月初六日，江苏巡抚杨魁题。}
\BookNote{30.}{中国第一历史档案馆，刑科题本，土地债务类，乾隆十九年十二月初二日，江宁巡抚庄有恭题。}
\BookNote{31.}{中国第一历史档案馆，刑科题本，婚姻家庭类，乾隆二十五年三月十九日，刑部尚书鄂弥达题。}
\BookNote{32.}{中国第一历史档案馆，刑科题本，婚姻家庭类，乾隆四十一年五月十九日，江宁巡抚杨魁题。}
\BookNote{33.}{中国第一历史档案馆，刑科题本，婚姻家庭类，乾隆四十一年五月十九日，江宁巡抚杨魁题。}
\BookNote{34.}{中国第一历史档案馆，刑科题本，土地债务类，嘉庆二十三年七月初七日，大学士管理刑部事务常煦题。}
\end{BookNotes}

% 原书第 27 页（PDF 第 40 页）
\BookPage{27}{40}
\noindent 刘氏供：丈夫吕文才开烧饼店生理。35 8.从事与「穿」有关职业者山西太谷人尤沅淳在扬州租房开张衣店，雇人管理账目36。这种衣店显然不是店主组织员工在店内制作衣服销售，而有专门的制衣商供应成衣。有人专门经营袜子店。

史二供：甘泉县人，二十一岁，父母在。小的住家与吕和友袜店相近，向来认识。乾隆四十七年五月十一日上午，小的到他店里买袜，吕和友拿一双夏布袜与小的看，说要三百文钱。小的把袜穿在脚上试看，觉得小些，要他另换。吕和友怪小的价未讲定把袜穿…。37还有人开洗衣服店为生，前面已有述及。9.各种小商贩这些小商贩多为流动经营者，但也多专销一种或以一种商品销售为主。刘大喜子供：甘泉县人，三十二岁，贩卖萝卜度日。38赵二受子供：江都县人，二十九岁，小的卖鱼生理。39钱二供：江都县人，二十九岁，平日卖烧饼度日。40孔兴名供：江都县人，小的哥孔兴广卖梨度日。41

\vfill
\begin{BookNotes}
\BookNote{35.}{中国第一历史档案馆，刑科题本，土地债务类，乾隆五十一年十一月十四日，江宁巡抚闵鹗元题。}
\BookNote{36.}{中国第一历史档案馆，刑科题本，土地债务类，道光元年五月十六日，江宁巡抚魏良煜题。}
\BookNote{37.}{中国第一历史档案馆，刑科题本，土地债务类，乾隆四十七年十一月初二日，江宁巡抚闵鹗元题。}
\BookNote{38.}{中国第一历史档案馆，刑科题本，土地债务类，嘉庆十四年六月十六日，大学士管理刑部事务董诰题。}
\BookNote{39.}{中国第一历史档案馆，刑科题本，土地债务类，乾隆五十七年七月十五日，江宁巡抚奇丰额题。}
\BookNote{40.}{中国第一历史档案馆，刑科题本，土地债务类，乾隆三十八年六月二十二日，大学士管理刑部事务刘统勋题。}
\BookNote{41.}{中国第一历史档案馆，刑科题本，土地债务类，道光六年十月二十七日，大学士管理刑部事务蒋攸铦题。}
\end{BookNotes}

% 原书第 28 页（PDF 第 41 页）
\BookPage{28}{41}
\BookSubheading{（五）开钱店为业}

\noindent 有关钱店的个案较少，我们见到一例个案中显示出这一行业的存在。宣廷魁在马克俊钱店帮伙。雍正十三年七月十七日，马克俊外出，王辅臣将银二钱四分向宣廷魁换钱，彼此较量钱数，引发口角42。

\BookSubheading{（六）与盐运有关的职业}

\noindent 周陈氏供：生有三子，从前都在盐船上相帮。43

\BookSubheading{（七）挑脚、赶脚和抬轿为业}

\noindent 扬州虽有行船运输之便，但陆上搬运仍不可少。挑脚也是当地的重要行业。其中赶脚为借助畜力运货。吴绍年之妻供：甘泉县人，丈夫赶脚营生，养有一驴。44还有专门挑运某一物品者。陈四供：江都县人，二十二岁，平日挑柴营生。45郭富供：甘泉县人，小的是抬戏箱营生，时常出门。46抬轿从业者。我们从多件个案中发现，扬州居民，特别是女性居民出行时雇轿很方便，说明此行业从业者不在少数。有的抬轿之人被富裕人家雇在家中，成为相对固定的轿夫。李德供：江都县人，五十一岁。乾隆七年上雇在王城家佣工，卞文也雇在家主家。乾隆十六年九月初一日，家主侄女往亲戚家去，是沈文同卞文抬轿送去，

\vfill
\begin{BookNotes}
\BookNote{42.}{中国第一历史档案馆，刑科题本，土地债务类，乾隆十四年六月十日，大学士兼管刑部事务来保题。}
\BookNote{43.}{中国第一历史档案馆，刑科题本，婚姻家庭类，乾隆三十一年八月初八日，江苏巡抚明德题。}
\BookNote{44.}{中国第一历史档案馆，刑科题本，土地债务类，乾隆五十八年二月二十二日，大学士管理刑部事务阿桂题。}
\BookNote{45.}{中国第一历史档案馆，刑科题本，土地债务类，乾隆十四年六月十九日，大学士兼管刑部事务来保题。}
\BookNote{46.}{中国第一历史档案馆，刑科题本，婚姻家庭类，乾隆元年七月二十九日，总理刑部事务允礼题。}
\end{BookNotes}

% 原书第 29 页（PDF 第 42 页）
\BookPage{29}{42}
晚上是小的同沈文抬轿接回，得有银封。初二日，小的把银封换钱三十三文，沈文抬过两次，分去十六文，卞文分得八文，小的分了九文。47

\BookSubheading{（八）以各种手艺为业}

严春源供：江都县人，年二十六岁。小的箍桶手艺。每日挑箍桶担在外做生意。48刘载刚供：江都县人，三十一岁，小的瓦匠生理。49

\BookSubheading{（九）以短工为业}

\noindent 短工是扬州人从事较多的一个行业，实际是打零工。他们多受雇与人，其雇主则可能遍布各种职业。

\BookSubheading{(十）乞丐}

\noindent 乞丐本非一种职业，但不少地方，特别是城市中，有专业乞丐。他们有的形成组织，俗称丐帮，有自己的活动地盘。扬州作为商业城市，即有这类以乞讨为生者。

王六供：江都县人，年三十九岁，父亲已故，母七十八岁，哥子王三、王五，兄弟王七，久已分居。小的与张大筛子、袁文儿、吴大、陈元珍、张乔山同在万寿州一带求乞度日。丐头宦贵每遇居民婚丧得钱，包揽看门，不容小的们求乞。平时见小的们讨得钱文，又要索诈酒食，稍不遂意就要打骂驱逐，心里忿恨。同案中宦余氏供：宦贵是小妇人丈夫，充当丐头。50

可见，扬州的乞丐的确具有专业性质，不仅求乞，而且有的包揽了居民婚丧活动的看门（维护秩序）之责。通过驱逐其他求乞者，从婚丧主人那里获得更多的利

\vfill
\begin{BookNotes}
\BookNote{47.}{中国第一历史档案馆，刑科题本，土地债务类，乾隆十八年三月初二日，署理刑部事务阿克敦题。}
\BookNote{48.}{《史料辑刊》，第一册，第307页。}
\BookNote{49.}{中国第一历史档案馆，刑科题本，婚姻家庭类，乾隆五十八年六月初九日，大学士管理刑部事务阿桂题。}
\BookNote{50.}{《史料辑刊》，第3册，第1691页。}
\end{BookNotes}

% 原书第 30 页（PDF 第 43 页）
\BookPage{30}{43}
\noindent 益。

当然，城中还有一些零散游丐。

\begin{BookQuote}
纪中子供：年十五岁，父亲年六十四岁，是双瞽，母亲已故，并无兄弟妻子。小的讨乞度日。51
\end{BookQuote}

在没有社会保障制度的时代，这些城市无业者，沦入以乞讨度日的境地是难免的。

\begin{BookQuote}
刘麻子供：甘泉县人，三十八岁，并无兄弟，也无妻室，在外求吃。52
\end{BookQuote}

值得注意的是，从个案可见，当时的雇佣劳动者往往需人推荐才能得到就业岗位。

\begin{BookQuote}
翟文斌供：二十八岁，与李三九认识，并无嫌隙。乾隆五十二年六月初，有龚大头因没生意托小的代荐帮工，适有开茶铺的李士旺曾托小的觅人帮伙，小的就把龚大头荐去，言明每月工钱三百文。六月十七日小的遇见李士旺，说龚大头已被李三九荐到李佩店内帮工，并埋怨小的荐人不妥。同案龚大头供：小的与翟文斌、李三九熟识。乾隆五十二年六月初小的因没生意，托翟文斌荐到李士旺茶馆内帮伙工，讲定每月工钱三百文。六月十六日，李三九来叫小的到李佩茶馆内暂帮几日。小的因李士旺尚没支过工钱，就应允
\end{BookQuote}

同去。53

从供词中可见，推荐人并没有从中获得什么好处。这种雇工方式实际带有一定保人性质。

综上所述，扬州人从事的职业有两大主要类别，一是与运河有关的「船运」之业，一是以住宿、饮食、杂货、生活用品修理等为主的商业服务业。它显示出扬州作为运河沿线商业消费城市和区域居民的职业特征。非农行业是扬州居民的主要从业领域。扬州城镇人口的吃、穿、住、行、用与市场建立起紧密关系，这说明清代

\vfill
\begin{BookNotes}
\BookNote{51.}{中国第一历史档案馆，刑科题本，土地债务类，嘉庆十九年二月初八日，江苏巡抚朱理题。}
\BookNote{52.}{中国第一历史档案馆，刑科题本，土地债务类，乾隆二十七年二月二十一日，刑部尚书舒赫德题。}
\BookNote{53.}{中国第一历史档案馆，刑科题本，土地债务类，乾隆五十三年八月初四日，刑部尚书喀宁阿题。}
\end{BookNotes}

% 原书第 31 页（PDF 第 44 页）
\BookPage{31}{44}
\noindent 中期扬州城市生活的品质是比较高的。

\BookSubheading{三、扬州的迁移、流动人口}

扬州作为清代中期重要的商业城市，不仅本地常住人口具有一定规模，而且汇集了不少流动人口。档案资料对此有一定表现。

\BookSubheading{（1）经商流动人口}

\begin{BookQuote}
崔百仓供：山东聊城人。嘉庆十七年十月，小的同姑夫姚文泰贩梨、枣来到扬州。因货未卖完，尚未回籍。54
\end{BookQuote}

该案背景资料显示，嘉庆十八年三月当事人仍在扬州。这意味着，他们在当地停留时间已超过五个月。该案例中所涉及的贩卖水果的山东人有五、六个，可见山东水果客商在当地成为一个群体。

\begin{BookQuote}
刘二供：小的是湖广武昌府江夏县人，三十一岁。乾隆二十四年冬小的相帮表兄周士贵撑船驾杉木来扬州，投程公来行里发卖。55杜六长供：江西南昌府新建县人，三十五岁。乾隆七年三月，小的与亲戚邓子文、邓配五三人各出六两银子合伙贩卖木篙到邵伯来卖，将木篙运到邵伯镇周东如行里发卖。六月十九日因木篙卖完了，二十日结算账目。56
\end{BookQuote}

以上三件个案有一个共同点是，外地客商携带货物来到扬州，需进入当地商行发卖，不能随意摆摊设点。我们认为，它说明两个问题，一是扬州当地城镇的货物销售是分区、分行业管理的，不同的货物，特别是相对大宗的货物集中在一起；二是外地客商进入扬州市场要通过当地的商行售卖，具有借地销售特征，而不允许擅自交易。另外，需要指出，外地客商入行销售期间，对自己的货物有完全支配权。即使已投入某一行，他们仍可将货物再运至其他行去销售。这说明，当地商人所开办的「行」并未垄断销售权，表现出有管理但又不失灵活的一面。

\vfill
\begin{BookNotes}
\BookNote{54.}{中国第一历史档案馆，刑科题本，土地债务类，嘉庆十八年十一月初一日，江宁巡抚朱理题。}
\BookNote{55.}{中国第一历史档案馆，刑科题本，土地债务类，乾隆二十五年十二月初三日，刑部尚书鄂弥达题。}
\BookNote{56.}{中国第一历史档案馆，刑科题本，土地债务类，乾隆八年五月二十二日，刑部尚书来保题。}
\end{BookNotes}

% 原书第 32 页（PDF 第 45 页）
\BookPage{32}{45}
\noindent 个案显示，外地人在扬州经营小生意者也有不少。周五供：丹徒县人，前往扬州钞关门外，做小生意度日。同案李胡子供：徐州府铜山县人，四十五岁。小的在钞关门外城河南岸租房居住，开杂货店生理，叫素识的范宗山在店帮伙。

同案范宗山供：河南固始县人，三十岁，小的到扬州遇见素识的李胡子，留小的在他店里帮伙。57孙六供：河南考城人，四十四岁，在邵伯河东开张饼店，并卖粉条生理。同案黄礼供：安徽砀山县人，四十四岁，在邵伯河西开张烧饼粉条店，雇一山东曹县人做帮工。孙六时常挪借货物的。58挪借货物实际为，当自己店中某种商品脱销时从邻居店中借来销售，以后将所借货物归还。

\BookSubheading{（2）各种流动性佣工}

\noindent 扬州的船运业所需雇工多数为外来人员。王世林供：山东东平州人，四十三岁。在江都拉度日。59邰子良供：盐城县人，今年三十五岁，向在扬州一带帮人拉船为生，没有船拉就在沈朝观饭店里吃饭住。小的一向因没有船拉，陆续欠沈朝观饭钱一两二钱三分。60蔡玉供：安徽凤阳人，二十七岁。帮人撑船度日。父故母在，并无兄弟妻子。小的在丁玉美船上雇做水手，丁玉美因无生意，陆续向小的借钱三千文，约俟装货归还…小的因家乡被水，要寄钱回，向丁玉美索欠…。61

其它行业的佣工也有不少外乡人。相对来说，流动择业群体中扬州附近州县人

\vfill
\begin{BookNotes}
\BookNote{57.}{中国第一历史档案馆，刑科题本，土地债务类，乾隆五十六年九月二十三日，江宁巡抚长麟题。}
\BookNote{58.}{中国第一历史档案馆，刑科题本，土地债务类，乾隆四十一年八月初八日，江宁巡抚杨魁题。}
\BookNote{59.}{中国第一历史档案馆，刑科题本，土地债务类，乾隆五十九年四月二十一日，江宁巡抚奇丰额题。}
\BookNote{60.}{中国第一历史档案馆，刑科题本，土地债务类，乾隆二十年一月二十日，署理两江总督尹继善题。}
\BookNote{61.}{中国第一历史档案馆，刑科题本，土地债务类，道光三十年三月初三日，江宁巡抚傅绳勋题。}
\end{BookNotes}

% 原书第 33 页（PDF 第 46 页）
\BookPage{33}{46}
\noindent 更多。吕美供：宝应县人，三十六岁，小的向在扬州挑水度日。62房有仁供：小的泰州人，四十岁，并无兄弟妻子，常在江都县地方帮做短工度日。乾隆五十一年七月小的到陈金书家借住寻工。63薛成锦供：小的丹徒县人，三十岁，在案下佣工度日。64

可见，扬州本地船运业、服务业和各种工商行业对雇佣劳动力的需求相对旺盛，因而它吸引大量周边地区成年劳动力前来谋生。

\BookSubheading{（3）农村无业和生存困难者的流入}

\noindent 一般来说，在农业与非农业、农村与城镇并存的时代，城镇常常会成为农村贫穷无业者

寻求新的谋生方式的去处。扬州在清代中期对周围农村人口具有这样的吸引力。

找寻工作、挣钱口是多数农村无业者来扬州的主要目的。

丁张氏供：盱眙县人。丈夫丁名因家里穷苦，没有生计，从前曾在扬州帮工。今年六月二十九日，带了小妇人并前夫生的两个儿子并小妇人生的周岁一个幼女，同到扬州来寻人家帮工度活。七月初九日到了扬州。65嵇太沅供：胡顺是小的表侄，他原籍为泰州。因父母俱故，别无亲属，到江都找寻工作，借住小的家内。66

也有来此地投亲或逃荒之人。

\vfill
\begin{BookNotes}
\BookNote{62.}{中国第一历史档案馆，刑科题本，婚姻家庭类，嘉庆十二年三月初三日，江宁巡抚汪日章题。}
\BookNote{63.}{中国第一历史档案馆，刑科题本，婚姻家庭类，乾隆五十二年五月十八日，江宁巡抚闵鹗元题。}
\BookNote{64.}{中国第一历史档案馆，刑科题本，土地债务类，嘉庆十三年五月初七日，大学士管理刑部事务董诰题。}
\BookNote{65.}{中国第一历史档案馆，刑科题本，土地债务类，乾隆三十闰二月十二日，大学士管理刑部事务刘统勋题。}
\BookNote{66.}{中国第一历史档案馆，刑科题本，土地债务类，乾隆五十九年四月二十一日，江苏巡抚奇丰额题。}
\end{BookNotes}

% 原书第 34 页（PDF 第 47 页）
\BookPage{34}{47}
王林氏供：小妇人丈夫王奇是海州人，本来有病，因穷苦难过来到扬州、同小妇人来投亲不遇，就租乔老汉房子暂住的。67彭大供：山东郯城人，二十岁。家里年荒，逃出来讨饭度日的。68

\BookSubheading{（4）流动人口向常住人口的转化}

\noindent 一般来说，城市的维持和稳定发展需要有一定数量的常住人口作为支撑，常住人口又多从流动人口转化而来。清代中期扬州保持了百年的繁盛期，它必然使一部分流动人口沉淀为常住人口。档案资料对此有所反映。符文德供：小的原籍浙江，来扬州已有两代，开张银铺。后因父母双亡，店就闭歇。今年七月，小的进京贸易，家内只有幼弟并妹子同小的妻子及老妈，雇仆吴二在家服役。69

当事人为浙江籍人，但即使父母去世、银铺闭歇，他和弟弟和妹妹也未返回故乡，而是以扬州为居家所在，出外从事贸易活动。无疑他们已成为本地常住人口。王徐氏供：王发来是小妇人丈夫，原籍山西太谷县，寄居江都。嘉庆二十五年三月丈夫同籍人尤沅淳来扬州，租赁小妇人家房屋开张衣店，叫丈夫在店管理账目，每月辛工钱两千文。70

王发来在当地已置有产业，有房屋出租给他人作店面，显然已不同于一般流动人口。

张二夸子供：三十二岁，原籍山东，寄居甘泉县。小的自幼来扬州，原籍何县，记不清了。父母俱故，并没有弟兄，娶妻戴氏，生有一女，赶猪营生。71

当事人小时候（应在嘉庆年间）随父母来扬州，与家乡已脱离联系，成为固定

\vfill
\begin{BookNotes}
\BookNote{67.}{中国第一历史档案馆，刑科题本，土地债务类，乾隆二十年六月初八日，江苏巡抚庄有恭题。}
\BookNote{68.}{中国第一历史档案馆，刑科题本，土地债务类，乾隆十三年十二月十五日，署理刑部尚书阿克敦题。}
\BookNote{69.}{中国第一历史档案馆，刑科题本，婚姻家庭类，乾隆十六年十一月初二日，江苏巡抚庄有恭题。}
\BookNote{70.}{中国第一历史档案馆，刑科题本，土地债务类，道光元年五月十六日，江宁巡抚魏元煜题。}
\BookNote{71.}{中国第一历史档案馆，刑科题本，土地债务类，道光二十年八月二十六日，署理江苏巡抚邵甲名题。}
\end{BookNotes}

% 原书第 35 页（PDF 第 48 页）
\BookPage{35}{48}
\noindent 居民。

可见，清代中期的扬州既是农村无业贫困人口寻求新生活的场所，也是对外地客商具有吸引力的经商和生活之地。

四结语

通过有限的档案资料，我们对清代中期扬州市镇经济发展水平和普通民众生活的一些方面加深了认识:扬州是一个商品经济发达、商业服务体系相对完善的城市和地区，表现为物流环节比较健全，商业活动形成规模，管理比较规范，商业经营的专业化程度较高。扬州作为非农业人口聚居之地，居民职业主要有两大主要类别，一是与运河有关的「船运」之业，一是以住宿、饮食、各类商品销售为主的商业服务业。它显示出扬州作为运河沿岸商业消费城市民众的职业特征。扬州还是东南地区流动人口的重要集散地，它既为农村无业贫困人口提供了一定的就业机会，也是对外地客商销售货物具有吸引力的城市。

不过，由于扬州的市镇中缺少具有一定规模的工场手工业，城市经济发展的基础并不坚实。一旦城市原有功能改变，消费人口数量减少，以商业服务业为主的市镇经济便会萎缩下去。


% 原书第 36 页（PDF 第 49 页）
\BookPage{36}{49}
